\section{보안 분석}
\label{sec:security}
\label{sec:analysis}

이 절은 세 질문을 다룬다. 정상 요청은 제대로 처리되는가,
틀린 거래나 결과를 성공으로 받아들일 수 있는가,
공개 원장을 보는 외부인이 숨긴 자료를 구별할 수 있는가이다.
이를 각각 정확성·건전성·프라이버시로 분석한다.

건전성에서는 참여자들이 공모해 잘못된 거래·증명·문서를 제출할 수
있다고 본다. 프라이버시에서는 정직하게 생성된 공개 이력을 읽는
외부 관찰자를 다룬다. 두 성질의 공격자 범위가 서로 다르다.
감사자는 신뢰된 주체로 두며, 프라이버시 보장은 비밀을 아는
거래 당사자나 승인된 감사자에게 적용하지 않는다.

결과는 다음 신뢰 가정 아래 성립한다. 원장의 권한 판정·현재 상태 검사·
원자적 갱신과 확정 이력의 인증 조회는 정확하고,
감사자와 상태 관리자는 신뢰된다고 둔다.
외부 절차는 유효한 감사키를 제공한다.
현실 입력의 진실성이나 키 전달 서비스 자체의 보안을 증명하는 것은 아니다.

암호 방식에는 증명의 근거 입력을 분석에서 확보할 수 있는 성질,
개방값의 결합성, 등록 증명의 건전성, 해시 충돌 저항성,
소유권 비밀의 유일한 표현과 허용 난수에서의 복호화 정확성을 요구한다.
분석의 공격자는 보안 매개변수에 대한 다항 시간·질의 수로 제한한다.
각 $\epsilon$은 해당 가정의 실패 가능성에 대한 전체 실행의 상계다.
정확한 공격 규칙·가정과 상세 증명은 부록~\ref{app:security-details}에 둔다.
정리는 이 가정들을 만족하는 구성에 적용된다.
실제 암호 설정·회로와의 대응은 \ref{sec:implementation-limits}절의
구현 검증 과제로 구분한다.

\subsection{정확성}
\label{sec:correctness-proof}

정확성은 규칙을 지킨 요청이 기대한 결과를 내는 성질이다.
사용 가능한 입력으로 올바른 분할 증명을 제출하면 두 출력을 기록하고,
정상 발급 문서의 공개 검증은 통과해야 한다.
필요한 기록과 유효 키가 있고 인증 조회가 끝나면,
고정 $H_t$의 추적도 정확한 이력과 미소비 집합을 반환해야 한다.

동결의 완료에는 추가 조건이 필요하다.
비어 있지 않은 전체 대상 $T$가 $H_t$부터 공통 $H_r$까지
소비되지 않고, 이미 차단된 대상과 새 동결의 차단 상태가 유지되어야 한다.
필요한 합법적 동결·조회·거래 결과와 공통 시점 확인도 정한 유한
한도 안에 확보해야 한다. 대상이 비면 빈 대상 결과를 반환한다.
경합이나 자료 부재가 있는 실행에 성공을 약속하는 정의는 아니다.

\textbf{정리 1 (Correctness).}
이 정상 조건 아래 정직 증명을 최대 $Q$회 생성하고,
한 번의 증명 완전성 실패 확률이 $\epsilon_{\rm comp}$ 이하이면
\begin{equation}
 \Pr[\mathsf{Fail}_{\rm corr}]
 \leq Q\epsilon_{\rm comp}+\epsilon_E+\epsilon_H.
 \label{eq:correctness-bound}
\end{equation}
이다. $\mathsf{Fail}_{\rm corr}$는 정상 조건이 성립하는 요청이
기대한 결과를 내지 못한 사건이다.
$\epsilon_E,\epsilon_H$는 전체 실행의 복호화 실패와 관련 해시
충돌의 상계다. 이는 정상 조건과 실패가 함께 생기는 확률이며,
정상 조건을 만족한 실행만 골랐을 때의 실패 비율은 아니다.

\emph{증명의 요지.}
올바른 증명은 검증을 통과하고, 유효한 실행 조건은 정해진 상태
갱신을 허용한다. 감사 암호문을 정확히 열면 실제 부모와 미래
식별자를 얻으므로 연결을 따라 정확히 추적한다.
동결의 추가 조건은 전체 대상의 공통 시점 차단을 보장한다.
따라서 실패는 증명 완전성·복호화·관련 해시 가정의 실패 중
하나를 필요로 한다. 세 가능성을 합하면 위 상계를 얻는다.
부록~\ref{app:correctness-details}는 정상 조건과 사건을 구체화한다.

\subsection{건전성}
\label{sec:soundness-proof}
\label{sec:binding-argument}

건전성은 틀린 거래나 결과를 올바른 것으로 받아들이기 어려운 성질이다.
여기서는 실제 생성 값과 다른 거래, 잘못된 추적 성공,
거짓 전체 완료와 차단 유지 중의 후속 소비를 하나의 실행에서 다룬다.
자료 부족의 명시적 실패나 대상을 보존한 \Alg{INCOMPLETE}는
공격 성공이 아니다.

\textbf{정리 2 (Soundness).}
부록~\ref{app:soundness-details}의 암호 가정과 위 신뢰 모델 아래
\begin{equation}
 \Pr[\mathsf{Exp}_{\rm sound}=1]\leq
 \epsilon_B:=\epsilon_X+\epsilon_C+\epsilon_M+\epsilon_H+\epsilon_E.
 \label{eq:bad-binding-bound}
\end{equation}
이다. $\epsilon_X,\epsilon_C,\epsilon_M,\epsilon_H,\epsilon_E$는
각각 증인 추출, 개방값 결합, 포함 관계, 관련 해시 충돌,
복호화의 전체 실행 실패 상계다.
증인 추출은 증명의 근거를 분석에서 확인하는 권한이며,
일반 검증자에게 비공개 자료를 공개한다는 뜻은 아니다.

\paragraph{거래와 추적의 일치.}
\label{sec:trace-argument}

증명의 첫 단계는 등록 때의 실제 값을 기준으로 거래를 확인하는 것이다.
수락된 증명의 근거 입력을 얻었을 때, 그것이 등록 때의 값과 다르면
개방값 결합성이 깨진다. 값과 적법한 소유 비밀이 같으면 소비
식별자도 같으므로, 현재 소비·상태 검사가 바로 그 객체에 적용된다.
감사 암호문도 같은 입력·출력에서 계산되므로 실제 부모와 미래
식별자를 돌려준다. 공개 Claim 검증의 문서·정책 연결도 같은 방식으로 다룬다.

다음으로 이 연결을 한 단계씩 따라간다.
실제 생성·소비 이력으로 정의한 그래프를 기준으로 삼으며,
감사 결과 자체를 정답으로 정의하지 않는다.
미래 식별자의 소비 기록이 0이면 정확히 미소비인 객체이고,
소비 기록이 있으면 실제 소비 이벤트의 전체 출력으로 이동한다.
이 과정을 반복하면 정해진 후속 경로와 미소비 집합을 얻는다.
선행 추적은 실제 부모를 거슬러 올라간다.
자료가 없거나 일관성이 깨지면 성공한 결과를 내지 않는다.
형식적인 시작점과 간선의 정의는 부록~\ref{app:trace-details}와 같다.

\paragraph{전체 차단의 보장.}

$T$를 $H_t$에서 찾은 전체 미소비 회계 증표·미해결 이전 증표로 두자.
공통 $H_r\geq H_t$에서 다음 조건이 성립한다고 하자.
\begin{equation}
 T\ne\varnothing,\qquad
 \forall t\in T:\ \B_{H_r}(t).
 \label{eq:frontier-cut}
\end{equation}
이는 대상이 비어 있지 않고 모두 미소비·차단 상태라는 뜻이다.
후속 사용을 계속하려면 남은 기록 중 적어도 하나를 먼저 써야 한다.
$H_t$ 이후 최초의 그런 사용이 있었다면 $T$의 객체 하나가 소비됐어야
하므로, $H_r$에서 모두 미소비라는 조건과 모순된다.
$H_r$ 이후에도 $T$의 차단이 유지되면 모든 소비 경로의 현재 상태
검사가 그 첫 사용을 거부한다. 다른 정상 입력과 합치는 거래도
차단된 입력의 검사를 피할 수 없다.

이 논증에 V절의 전체 대상·공통 시점·거래 결과 확인을 적용하면
거짓 전체 완료도 배제된다. 전이·추적·완료의 실패는 같은 암호
가정의 실패 집합에 포함되므로 상계를 세 번 더하지 않는다.
완료 보고와 오류가 함께 발생하는 확률의 상계이며,
완료를 보고한 실행만 고른 조건부 오류율은 아니다.

V절의 동결·해제 사례는 공통 시점 없이 완료를 보고하면 안 되는 이유다.
동결 전에 소비된 대상을 목록에서 빼는 것도 남은 파생 기록을 놓치게 한다.
반대로 \Alg{INCOMPLETE}라고 해서 현재 쓸 수 있는 후속 객체가
반드시 있다는 뜻은 아니다. 대상이 종료로 소진됐거나 다른 실행이
파생 객체를 이미 막았을 수도 있다.
이 결과는 본 절차의 충분조건이며 모든 차단 방식의 필요조건은 아니다.

\subsection{프라이버시}
\label{sec:privacy-proof}

프라이버시는 숨긴 수치나 입력 연결은 다르지만,
공개하기로 한 정보는 같은 두 이력을 외부인이 구별하기 어려운 성질이다.
예를 들어 거래 종류와 공개 제품 주장이 같은 두 유효한 이력에서,
실제로 쓴 원료 기록의 연결이나 비공개 수치가 다를 수 있다.
다른 공개 조건도 같다면 외부인이 어느 이력인지 알아내기 어려워야 한다.
거래 종류·정책·입출력 형상·호출자·시간·기한·공개 Claim과
공개 색인의 동일성 같은 정보는 보호 대상에서 제외한다.
정확한 공개 범위는 부록~\ref{app:privacy-details}에 명시한다.

대상은 독립적인 정직 난수로 생성된 이력을 읽는 수동 외부 관찰자다.
설정 전에 고정한 $K-1$개 이하 위원회 조각은 볼 수 있지만,
참여자 비밀·개방값, $SK_B$, 보호 관계의 감사·상태 공개는 받지 않는다.
능동 거래 제출·복호화 질의도 제외한다.
따라서 승인된 감사자나 비밀을 아는 거래 당사자까지 보호하는 주장은 아니다.

분석에서는 용도가 다른 해시를 독립된 이상적 무작위 함수로 다루는
무작위 오라클 모형을 사용한다.
두 키 암호화의 기밀성과 여러 증명을 함께 보아도 성립하는 영지식
성질도 가정한다.

\textbf{정리 3 (Privacy).}
부록~\ref{app:privacy-details}의 게임과 충분조건 아래
\begin{equation}
 \mathsf{Adv}_{\rm hist}\leq
 2(\epsilon_Z+\epsilon_A+\epsilon_D+\epsilon_{\rm RO}).
 \label{eq:privacy-bound}
\end{equation}
이다. $\epsilon_Z$는 전체 증명의 영지식 구별 이득,
$\epsilon_A,\epsilon_D$는 임계 키 보호와 데이터 암호화의 구별 이득이다.
$\epsilon_{\rm RO}$는 숨은 입력 추측·관련 충돌의 상계다.
각 이득은 두 실험을 구분할 수 있는 정도를 나타낸다.

\emph{증명의 요지.}
공개 증명을 비공개 입력 없이 만들 수 있는 분석용 자료로 바꾸고,
키 보호·감사 암호문의 내용을 대체한 뒤,
숨은 기록과 식별자를 독립된 무작위 값으로 바꾼다.
각 변화를 구별할 수 있으면 해당 가정도 깨뜨릴 수 있다.
마지막에 남는 분포는 공개 정보에만 의존하므로 두 이력은 구별하기 어렵다.
독립 키, 비밀·난수의 범위, 공개 자료의 재현과 각 변경 순서에 필요한
조건은 부록에서 확인한다.
